\documentclass[12pt, a4paper]{article}

% --- PAQUETES BÁSICOS ---
\usepackage[utf8]{inputenc}
\usepackage[spanish]{babel}
\usepackage[margin=2.5cm]{geometry}
\usepackage{graphicx}
\usepackage{amsmath}
\usepackage{hyperref}
\usepackage{tabularx} % Añadido para tablas auto-ajustables
\usepackage{array}    % Añadido para mejor formato en columnas

% --- CONFIGURACIÓN DE HIPERVÍNCULOS ---
\hypersetup{
    colorlinks=true,
    linkcolor=blue,
    filecolor=magenta,      
    urlcolor=cyan,
    pdftitle={Resumen ESP32 y ESP32S3},
    pdfpagemode=FullScreen,
}

\title{\textbf{Resumen uso ESP32 y ESP32S3}}
\author{Jairo David Diaz Luna}
\date{\today}

\begin{document}

\maketitle
\thispagestyle{empty} 
\newpage

\pagenumbering{roman} 
\tableofcontents
\newpage

\pagenumbering{arabic} 

\section{Introducción}

\section{Pines}

\section{Procesador}
\subsection{Caracteristucas}
\subsubsection{Accumulator Register ACCX}
Consta de un banco de 16 registros, util para operaciones MAC de matrices en memoria y ejecutarlas eficientemente ademas redondeando floats evitando desbordamientos 
dado a que posee registros de 40 bits con 4 canales , permitiendo que sea mas complejo desbordar realizando multiplicaciones iterativamente.
Finalmente realiza un bit-shift al final y reduce la salida a los tipos de variable manejados por la CPU.
\paragraph{Multiply Accumulate MAC}
Calcula un producto de dos numeros y lo suma a un acumulador.
Recordando la estructura de un float32: 
Signo (1 bit): Determina si es positivo o negativo.
Exponente (8 bits): Define la magnitud o escala del número (cuántas posiciones se mueve el punto decimal).
Mantisa o Fracción (23 bits): Almacena los dígitos significativos reales. El hardware asume un "1" implícito a la izquierda, por lo que internamente el procesador opera con 24 bits de precisión.

El algoritmo de multiplicacion consta en: 
\begin{enumerate}
    \item El hardware toma $A$ (16 bits) y $B$ (16 bits). Los multiplica. El resultado exacto es $C$ (32 bits). Ese $C$ se guarda en el acumulador de 40 bits.
    \item El hardware toma datos nuevos $D$ (16 bits) y $E$ (16 bits). Los multiplica. El resultado es un nuevo valor $F$ (32 bits).
    \item El hardware toma ese nuevo $F$ y se lo suma a lo que ya estaba en el acumulador ($C$).
\end{enumerate}
Por cada ciclo de acumulacion por acarreo se incrementa un bit, implica que $2^n$ multiplicaciones requieren solo $n$ bits de acarreo. 

\subsubsection{Quantized Accumulator QAAC (ESP32S3)}
Exclusivo de la ESP32S3, consta del ACCX mencionado pero con cuatro canales, siendo cuadruple. Util para modelos de matrices iterativas, como redes neuronales multicapa o metodos iterativos de ecuaciones. 

\subsubsection{Arithmetic Logic Unit}
Bloque combinacional que ejecuta tanto multiplicaciones, adiciones, operaciones logicas AND, OR y NOT y registros QR.

\subsubsection{Registros}
Los Xtensa poseen varios registros, tanto generales como vectoriales.
Los generales son
\begin{enumerate}
    \item AR usado para proposito general siendo 16 regsitros de 32 bits.
    \item FR usado solo para float32 siendo 16 registros. 
\end{enumerate}
Los vectoriales son
\begin{enumerate}
    \item QR usado para proposito general en paralelo siendo 8 registros de 128 bits
    \item SAR usado para bitshifting siendo 1 de 6 bits.
    \item ACCX y QAAC los usados de acumuladores de operaciones MAC.
    \item FFT usado para operaciones de FFT. 
    \item UA_State usado para guardar direcciones siendo un registro de 128 bits
\end{enumerate}

\subsection{Ultra Low Power Processor ULP}
Se encarga de las tareas de bajo consumo y de despertar al procesador principal. 
Este posee 8Kb de memoria RTC, reloj lento de 17.5 MHz.
Permite controlar perifericos como los dos ADC internos que posee, protocolo I2C, gestion de puertos Input/Output y sensores internos pasivos como temperatura o las entradas de modulos touch.

\section{XDXD}

\section{XDXD}

\section{XDXD}

\end{document}